\chapter{Terms, Substitutions, Matching, and Unification}
\label{ch:term-matching-unification}

A rewrite such as $\operatorname{add}(x,0)\to x$ is not an instruction to replace every source variable named x.
Here $x$ is a metavariable ranging over entire subterms.
Applying the rule requires finding one consistent substitution, instantiating the replacement, and changing exactly the selected occurrence.
The distinction between matching and solving equations determines both correctness and implementation cost.

\section{Give terms a signature and positions}
\label{sec:ch82-signature}
A first-order signature assigns an arity to each function symbol.
Constants have arity zero; variables belong to a separate set $X$.
Terms are variables or expressions $f(t_1,\ldots,t_n)$ where the signature gives $f$ arity $n$.
The companion represents variables with a distinct \code{Var} class and applications with immutable tuples, so the constant named x is not accidentally a metavariable.
Its small trusted examples supply well-formed signatures; a user-facing rule loader would additionally validate a global arity table.

A \term{term position} is a sequence of child indices; the empty sequence selects the root.
The code uses zero-based indices, so position $(0,1)$ selects the second child of the first child.
We write $t|_p$ for the subterm at $p$ and $t[u]_p$ for replacing that occurrence with $u$.
Figure~\ref{fig:ch82-position} selects the constant one in $\operatorname{mul}(\operatorname{add}(v,0),1)$.

\begin{figure}[htbp]
\centering
\begin{tikzpicture}[node distance=13mm and 14mm]
\node[basic block] (m) {mul: root};
\node[basic block,below left=of m] (a) {add: (0)};\node[basic block,below right=of m] (one) {one: (1)};
\node[basic block,below left=of a] (v) {v: (0,0)};\node[basic block,below right=of a] (z) {zero: (0,1)};
\draw[cfg edge] (m)--(a);\draw[cfg edge] (m)--(one);\draw[cfg edge] (a)--(v);\draw[cfg edge] (a)--(z);
\end{tikzpicture}
\caption{A child-index path identifies one occurrence independently of its printed spelling.}
\label{fig:ch82-position}
\end{figure}

A rule $l\to r$ requires a nonvariable left side and $\operatorname{Vars}(r)\subseteq\operatorname{Vars}(l)$.
The latter condition prevents an ordinary first-order step from inventing unbound metavariables.
Fresh-name generation and binding-aware rewriting are useful extensions, but need separate semantics rather than silently violating that rule contract.
The lab rejects these malformed rules before applying them.

\section{Substitute consistently}
\label{sec:ch82-substitution}
A substitution maps finitely many variables to terms and extends homomorphically through application nodes.
Repeated variables denote repeated occurrences of the same substituted term.
For $\sigma=\{x\mapsto a,y\mapsto g(a)\}$, applying $\sigma$ to $f(x,y)$ gives $f(a,g(a))$.
Substitution composition order matters: applying one map after another is not generally a commutative dictionary union.

The companion recursively resolves chained substitutions and returns an idempotent solved map from unification.
It does not accept cyclic substitution maps supplied by arbitrary clients.
Its occurs check prevents the solver from constructing one.
For graph-based types that deliberately admit rational infinite trees, omitting the occurs check would describe a different mathematical domain, not merely a faster implementation of finite-term unification.

\section{Match patterns in one direction}
\label{sec:ch82-match}
\term{Matching} solves $l\sigma=t$ by instantiating variables of the pattern $l$, leaving the subject fixed.
For pattern $\operatorname{pair}(x,x)$, subject $\operatorname{pair}(a,a)$ succeeds with $x\mapsto a$.
Subject $\operatorname{pair}(a,b)$ fails for distinct constants.
Binding x to a at its first occurrence and overwriting it with b at the second is an incorrect matcher.
Table~\ref{tab:ch82-match} makes the repeated-variable check explicit.

\begin{table}[htbp]
\caption{Repeated pattern variables impose equality constraints on the subject.}
\label{tab:ch82-match}
\centering
\begin{tabular}{lll}
\toprule
Pattern & Subject & Result \\
\midrule
pair(x,x) & pair(a,a) & x maps to a \\
pair(x,x) & pair(a,b) & failure \\
pair(x,y) & pair(a,b) & x maps to a; y maps to b \\
f(x) & g(a) & root-symbol failure \\
\bottomrule
\end{tabular}
\end{table}

Algorithm~\ref{alg:ch82-match} uses a worklist of pattern/subject pairs.
With constant-time binding-table operations and constant-time equality of shared subject identities, matching is linear in visited pattern structure.
The lab uses structural equality for repeated bindings, so comparisons of large repeated subterms add their actual traversal costs.
Hash-consed nodes or cached structural hashes can accelerate those checks, but hash agreement alone is not exact equality unless collisions are checked.

\begin{algorithm}[htbp]
\caption{Match a finite first-order pattern against one fixed subject.}
\label{alg:ch82-match}
\begin{pseudocode}
\Require well-formed pattern and subject; separate pattern-variable namespace
\Ensure consistent matcher substitution or failure
\State initialize pair worklist and empty bindings
\While{a pair $(p,s)$ remains}
\If{$p$ is a pattern variable}
\State bind it to $s$ if unbound; otherwise require equality with its binding
\Else
\State require matching root symbol and arity
\State enqueue corresponding children
\EndIf
\EndWhile
\State return bindings
\end{pseudocode}
\end{algorithm}

\section{Unify in both directions with an occurs check}
\label{sec:ch82-unify}
\term{Unification} solves equations by instantiating variables on either side.
For $f(x,g(y))=f(g(a),g(a))$, decomposition gives $x=g(a)$ and $y=a$.
The solved substitution makes both original terms identical.
For $x=f(x)$ there is no finite-term solution because any candidate would contain itself as a proper subterm.

Algorithm~\ref{alg:ch82-unify} dereferences existing bindings before every equation decision.
The occurs check examines the dereferenced right side, not merely its original syntax, so an indirect cycle through earlier bindings is also rejected.
An equation $x=x$ is discarded before the occurs check; it is not an error.

\begin{algorithm}[htbp]
\caption{Unify finite first-order terms using an explicit equation worklist.}
\label{alg:ch82-unify}
\begin{pseudocode}
\Require finite terms with disjoint symbol and variable representations
\Ensure most-general finite-term unifier, or failure
\State initialize worklist with the input equation; substitution empty
\While{an equation remains}
\State dereference both sides through the current substitution
\If{sides are equal} \State discard equation
\ElsIf{one side is a variable $x$}
\State orient variable left; reject if $x$ occurs in the other side
\State record its binding
\Else
\State require equal root symbols and arities; enqueue child equations
\EndIf
\EndWhile
\State resolve substitution chains before returning the solved map
\end{pseudocode}
\end{algorithm}

This elementary algorithm is described for clarity, not as a linear-time unification implementation.
Repeated substitution expansion, occurs traversals, and copying can be expensive; explicitly materializing a shared solved substitution as trees can even require exponential output.
Graph-based algorithms with union-find and carefully maintained occurs information offer stronger bounds under their own assumptions.
The chapter's type-inference extension uses the same finite-type occurs principle while keeping its solver contract separate.

\section{Replace occurrences and replay evidence}
\label{sec:ch82-rewrite}
A rewrite step selects a position $p$, matches $l$ against $t|_p$, and constructs $t[r\sigma]_p$.
Immutable tree rebuilding copies only the path's ancestors if child subtrees can be shared safely.
The lab records rule index, position, and matcher substitution and can replay the trace independently of the search strategy.
Figure~\ref{fig:ch82-trace} separates discovery from verification.

\begin{figure}[htbp]
\centering
\begin{tikzpicture}[node distance=13mm]
\node[basic block] (search) {search positions and rules};
\node[basic block,below=of search] (event) {rule ID, position, substitution};
\node[basic block,below=of event] (replay) {recheck match; rebuild replacement};
\draw[cfg edge] (search)--(event);\draw[cfg edge] (event)--(replay);
\end{tikzpicture}
\caption{Replay validates the claimed local step without repeating redex discovery.}
\label{fig:ch82-trace}
\end{figure}

For $\operatorname{mul}(\operatorname{add}(v,0),1)$, outermost-first uses the multiplication identity at the root and then the addition identity at the root of the result.
Innermost-first instead rewrites position $(0)$ and then the root.
Both reach v in two steps for this rule set, but their certificates differ.
Table~\ref{tab:ch82-trace} records the selected paths.

\begin{table}[htbp]
\caption{Two strategies can agree on a result while selecting different redex positions.}
\label{tab:ch82-trace}
\centering
\begin{tabular}{lll}
\toprule
Strategy & First position/rule & Second position/rule \\
\midrule
outermost-first & root, multiply by one & root, add zero \\
innermost-first & child zero, add zero & root, multiply by one \\
\bottomrule
\end{tabular}
\end{table}

These are first-order terms without object-language binding.
Rewriting lambda bodies requires alpha-equivalence and capture avoidance or an explicit binder representation such as declaration IDs or de Bruijn indices.
Treating the lambda parameter as an ordinary child and substituting by spelling can capture a free variable.
Chapter~\ref{ch:functional-core-semantics} introduces the environment-based alternative used by the functional implementation.

\section{Explain why the elementary unifier is most general}
\label{sec:ch82-generality}
Unification does more than find a convenient assignment.
A \term{most-general unifier} represents all solutions by leaving unconstrained variables free: any other unifier factors through it by a further substitution.
The elementary equation transformations preserve precisely that solution set.
Discarding an equation t=t removes no constraint; decomposing equal-root applications replaces one equation by exactly the child equalities needed to satisfy it.
Different root symbols or arities cannot become equal by substituting variables below those roots.

Eliminating x=t, when x does not occur in t, commits to an equality that every solution already has to satisfy.
Replacing later uses of x by t expresses the remaining equations in the smaller unknown space.
The occurs check is what keeps this elimination within finite terms.
After the worklist is empty, every original equation is satisfied by the accumulated substitution; variables not forced to a concrete structure remain available for arbitrary further instantiation.
That is the reason for generality, not a consequence of picking the first constant the implementation sees.

For equation pair(x,y)=pair(y,z), one solved form maps x to z and y to z while leaving z free.
Mapping all three to a happens to solve the equation, but is less general: it cannot represent the solution mapping all three to b unless a itself were a metavariable rather than a fixed constant.
The distinction matters for type inference, where prematurely choosing Int for an unconstrained unknown destroys polymorphism.
The lab verifies every returned substitution by applying it to both original terms, while this solution-preservation argument explains the stronger generality contract.

\section{Keep matching modulo equations and binding outside the basic contract}
\label{sec:ch82-modulo}
First-order syntactic matching compares ordered children under exact symbols.
Matching modulo commutativity would allow add(a,b) to match add(x,a) by swapping children, and associativity can require choosing among several groupings.
Flattening and sorting can provide a canonical representation for a deliberately associative/commutative operator, but only after its arithmetic and effect semantics permit those equations.
The ordinary ordered-child matcher does not silently acquire those capabilities.

Higher-order patterns add another dimension because metavariables may denote functions and terms contain binders.
A replacement can capture a formerly free variable unless the representation tracks binding identities or performs capture-avoiding substitution.
An environment evaluator avoids repeated source substitution during execution, but a compiler transformation that rewrites lambda syntax still needs its own binding-safe representation.
The functional optimizer in this expansion restricts itself to local structural identities and traverses binder forms explicitly; it does not advertise a general higher-order unifier.

Conditional rewriting also separates structural substitution from premise solving.
A pattern can match perfectly while an arithmetic range, purity, or dominance condition is unproved.
Such a match is a candidate, not an admitted transformation.
The matching substitution should remain immutable while semantic conditions are checked, and a graph edit that invalidates a condition must cause revalidation before the replacement is committed.
This is why the rule-engine chapter treats matching, analysis queries, and commit as distinct stages.

\section{Exercises}
\label{sec:ch82-exercises}
\begin{enumerate}
\item\label{ex:ch82-match} Match pair(x,x) against pair(a,a) and pair(a,b), showing every binding check.
\item\label{ex:ch82-unify} Unify $f(x,g(y))$ and $f(g(a),g(a))$ and verify the result by substitution.
\item\label{ex:ch82-occurs} Solve $x=y$ followed by $y=f(x)$ and identify the dereferenced occurs-check failure.
\item\label{ex:ch82-position} In the worked multiply/add term, give the position of zero and the result of replacing only that occurrence with one.
\end{enumerate}

\section{Further reading}
\label{sec:ch82-reading}
Ohlebusch supplies the first-order rewrite and unification framework used here~\cite{ohlebusch-2002-rewriting}.
Appel connects substitution-like reasoning to compiler representations and environments~\cite{appel-1998-modern}.
